\documentclass[12pt,a4paper]{article}

% ── Packages ──────────────────────────────────────────────────────────────────
\usepackage[margin=2.5cm]{geometry}
\usepackage{amsmath}
\usepackage{graphicx}
\usepackage{booktabs}
\usepackage{hyperref}
\usepackage{caption}
\usepackage{subcaption}
\usepackage{float}
\usepackage{enumitem}
\usepackage{parskip}
\usepackage{xcolor}
\usepackage{microtype}

\hypersetup{
  colorlinks=true,
  linkcolor=blue!70!black,
  urlcolor=blue!70!black,
  citecolor=blue!70!black,
}

\newcommand{\placeholder}[1]{\textcolor{red!70!black}{\textit{[#1]}}}

% ── Title ─────────────────────────────────────────────────────────────────────
\title{%
  \textbf{Data Analysis in Psycholinguistics}\\[0.5em]
  \large Corpus and Experimental Data — Assignment Report
}
\author{%
  \placeholder{Your Name}\\
  \placeholder{Student ID}\\[0.4em]
  \small\textit{Linguistic Data 3: Data Modelling in Indian Languages}
}
\date{\today}

% ══════════════════════════════════════════════════════════════════════════════
\begin{document}
\maketitle
\tableofcontents
\newpage

% ══════════════════════════════════════════════════════════════════════════════
\section{Corpus Data Analysis: Hindi and Telugu Treebanks}
% ══════════════════════════════════════════════════════════════════════════════

\subsection*{Data and Tools}
The analysis uses the Telugu dependency treebank
(\url{https://github.com/ltrc/telugu_treebank}) and the Hindi Dependency
Treebank (HDTB pre-release v0.05,
\url{https://ltrc.iiit.ac.in/treebank_H2014/}).  Both treebanks are in
CoNLL-U / CoNLL-2006 format.  Parsing and analysis were performed in Python~3
using \texttt{pandas}, \texttt{scipy}, \texttt{matplotlib}, and
\texttt{seaborn}.  A custom parser handles both 8-column (CoNLL-2006) and
10-column (CoNLL-U) formats and skips malformed lines.

% ──────────────────────────────────────────────────────────────────────────────
\subsection{Q1: Distribution of Dependency Distances}
% ──────────────────────────────────────────────────────────────────────────────

The dependency distance for a token is defined as $|i - h_i|$, where $i$ is
the token index and $h_i$ is its head index.  Root tokens ($h_i = 0$) are
excluded.

\begin{table}[H]
\centering
\caption{Dependency distance summary statistics.}
\begin{tabular}{lrrrr}
\toprule
\textbf{Language} & \textbf{Mean} & \textbf{Median} & \textbf{Std} & \textbf{N} \\
\midrule
Telugu & \placeholder{X.XX} & \placeholder{X} & \placeholder{X.XX} & \placeholder{N} \\
Hindi  & \placeholder{X.XX} & \placeholder{X} & \placeholder{X.XX} & \placeholder{N} \\
\bottomrule
\end{tabular}
\label{tab:dist_stats}
\end{table}

\begin{figure}[H]
  \centering
  \begin{subfigure}[b]{0.48\textwidth}
    \includegraphics[width=\textwidth]{results/dependency_distance_histogram.png}
    \caption{Density histogram.}
  \end{subfigure}
  \hfill
  \begin{subfigure}[b]{0.48\textwidth}
    \includegraphics[width=\textwidth]{results/dependency_boxplot.png}
    \caption{Boxplot.}
  \end{subfigure}
  \caption{Dependency distance distributions for Telugu (blue) and Hindi (red).}
  \label{fig:dist}
\end{figure}

Both distributions are right-skewed: most dependencies are short-range (distance
1--3), with a long tail of long-distance dependencies.  The median is more
representative than the mean given this skewness.

% ──────────────────────────────────────────────────────────────────────────────
\subsection{Q2: Top-10 Dependency Relations}
% ──────────────────────────────────────────────────────────────────────────────

\begin{table}[H]
\centering
\caption{Top 10 dependency relations (raw counts) for Telugu and Hindi.}
\begin{tabular}{lrlr}
\toprule
\multicolumn{2}{c}{\textbf{Telugu}} & \multicolumn{2}{c}{\textbf{Hindi}} \\
\cmidrule(r){1-2}\cmidrule(l){3-4}
\textbf{Relation} & \textbf{Count} & \textbf{Relation} & \textbf{Count} \\
\midrule
\placeholder{r1} & \placeholder{n} & \placeholder{r1} & \placeholder{n} \\
\placeholder{r2} & \placeholder{n} & \placeholder{r2} & \placeholder{n} \\
\placeholder{r3} & \placeholder{n} & \placeholder{r3} & \placeholder{n} \\
\placeholder{r4} & \placeholder{n} & \placeholder{r4} & \placeholder{n} \\
\placeholder{r5} & \placeholder{n} & \placeholder{r5} & \placeholder{n} \\
\placeholder{r6} & \placeholder{n} & \placeholder{r6} & \placeholder{n} \\
\placeholder{r7} & \placeholder{n} & \placeholder{r7} & \placeholder{n} \\
\placeholder{r8} & \placeholder{n} & \placeholder{r8} & \placeholder{n} \\
\placeholder{r9} & \placeholder{n} & \placeholder{r9} & \placeholder{n} \\
\placeholder{r10} & \placeholder{n} & \placeholder{r10} & \placeholder{n} \\
\bottomrule
\end{tabular}
\label{tab:deprel}
\end{table}

\begin{figure}[H]
  \centering
  \includegraphics[width=0.85\textwidth]{results/dependency_relations_comparison.png}
  \caption{Grouped bar chart comparing the top dependency relations in Hindi and Telugu.}
  \label{fig:deprel}
\end{figure}

% ──────────────────────────────────────────────────────────────────────────────
\subsection{Q3: Notable Similarities and Differences in Syntactic Relations}
% ──────────────────────────────────────────────────────────────────────────────

\textbf{Similarities.}
Both Hindi and Telugu are verb-final (SOV) languages, and this is reflected in
their dependency profiles.  Nominal modifier (\texttt{nmod}) and noun-phrase
relations appear prominently in both treebanks.  Relations associated with
coordination (\texttt{cc}, \texttt{conj}) are present in both at comparable
rates, suggesting that coordinate structure is a shared strategy regardless of
language family.

\textbf{Differences.}
\begin{itemize}[leftmargin=1.5em]
  \item \textbf{Subject/object marking.}  Telugu, as a Dravidian language, uses
        case-marked postpositional phrases and has a relatively free word order
        within clauses.  Its subject dependency (\texttt{nsubj}) and object
        dependency (\texttt{obj} / \texttt{dobj}) counts therefore span a wider
        range of positions, potentially yielding longer average distances.  Hindi,
        while also relatively free, has stronger verb-agreement morphology that
        anchors subjects closer to the predicate.
  \item \textbf{Postpositional vs.\ adpositional relations.}  Telugu shows a
        higher rate of postpositional relations (e.g.\ \texttt{nmod:poss},
        postpositional case-markers) because all adpositions follow their
        nominal heads.
  \item \textbf{Light-verb constructions.}  Hindi exhibits a high frequency of
        complex predicates (e.g.\ $V + \textit{lenā/denā}$), typically annotated
        as \texttt{compound:prt} or similar.  This relation may appear
        disproportionately in the Hindi top-10 but not in Telugu.
\end{itemize}

% ──────────────────────────────────────────────────────────────────────────────
\subsection{Q4: Dependency Distance and Word Order Typology}
% ──────────────────────────────────────────────────────────────────────────────

\subsubsection*{Descriptive Comparison}

Table~\ref{tab:dist_stats} shows that \placeholder{[Telugu/Hindi]} has a
higher mean dependency distance (\placeholder{X.XX} vs.\ \placeholder{X.XX}).
Both languages are SOV, which places heavy NPs and embedded clauses before the
verb, creating long-distance dependencies between arguments and their predicates.

\subsubsection*{Significance Testing}

Because dependency distances are not normally distributed (right-skewed, count
data), a non-parametric \textbf{Mann-Whitney U test} was used.

\begin{table}[H]
\centering
\caption{Mann-Whitney U test results for dependency distances.}
\begin{tabular}{lr}
\toprule
\textbf{Statistic} & \textbf{Value} \\
\midrule
U statistic & \placeholder{U} \\
$p$-value   & \placeholder{p} \\
Significant ($\alpha = 0.05$)? & \placeholder{Yes/No} \\
Cohen's $d$ (effect size) & \placeholder{d} \\
\bottomrule
\end{tabular}
\label{tab:signif}
\end{table}

\textbf{Interpretation.}
\placeholder{If significant:} The difference in dependency distances is
statistically significant ($p < 0.05$), indicating that the two languages do
not share the same distance distribution.  A Cohen's $d$ of \placeholder{d}
suggests a \placeholder{small/medium/large} effect.  The longer mean distance
in \placeholder{language} may reflect \placeholder{freer postpositional phrase
ordering / heavier NP fronting / a richer system of case-marked modifiers}.
Both languages' SOV typology, however, constrains the overall shape of the
distribution to be right-skewed with a mode at distance~1.

% ──────────────────────────────────────────────────────────────────────────────
\subsection{Q5: Morphological Features}
% ──────────────────────────────────────────────────────────────────────────────

Morphological features were extracted from the CoNLL-U \texttt{feats} column
(format: \texttt{Feature=Value}, separated by \texttt{|}). Feature
\emph{category} counts are shown below.

\begin{table}[H]
\centering
\caption{Most frequent morphological feature categories in each language.}
\begin{tabular}{lrlr}
\toprule
\multicolumn{2}{c}{\textbf{Telugu}} & \multicolumn{2}{c}{\textbf{Hindi}} \\
\cmidrule(r){1-2}\cmidrule(l){3-4}
\textbf{Feature} & \textbf{Count} & \textbf{Feature} & \textbf{Count} \\
\midrule
\placeholder{f1} & \placeholder{n} & \placeholder{f1} & \placeholder{n} \\
\placeholder{f2} & \placeholder{n} & \placeholder{f2} & \placeholder{n} \\
\placeholder{f3} & \placeholder{n} & \placeholder{f3} & \placeholder{n} \\
\placeholder{f4} & \placeholder{n} & \placeholder{f4} & \placeholder{n} \\
\placeholder{f5} & \placeholder{n} & \placeholder{f5} & \placeholder{n} \\
\bottomrule
\end{tabular}
\label{tab:morph}
\end{table}

\begin{figure}[H]
  \centering
  \begin{subfigure}[b]{0.48\textwidth}
    \includegraphics[width=\textwidth]{results/telugu_morph_features.png}
    \caption{Telugu.}
  \end{subfigure}
  \hfill
  \begin{subfigure}[b]{0.48\textwidth}
    \includegraphics[width=\textwidth]{results/hindi_morph_features.png}
    \caption{Hindi.}
  \end{subfigure}
  \caption{Top morphological feature categories by token count.}
  \label{fig:morph}
\end{figure}

\begin{figure}[H]
  \centering
  \includegraphics[width=0.75\textwidth]{results/morph_feature_heatmap.png}
  \caption{Normalised morphological feature heatmap (\% of total feature tokens).}
  \label{fig:heatmap}
\end{figure}

% ──────────────────────────────────────────────────────────────────────────────
\subsection{Q6: Morphological Patterns and Typological Distinctions}
% ──────────────────────────────────────────────────────────────────────────────

\textbf{Indo-Aryan (Hindi) vs.\ Dravidian (Telugu).}

Hindi, belonging to the Indo-Aryan branch, retains a two-gender (masculine /
feminine) system inherited from Sanskrit via Prakrits, and marks agreement on
verbs for gender, number, and (partially) person.  Consequently, the
\texttt{Gender} feature is expected to appear at high frequency in the Hindi
treebank, attaching to nouns, adjectives, and inflected verbs alike.  Case is
marked by postpositions and partially by nominal inflection, with the direct
and oblique distinction being the primary nominal case contrast.

Telugu, as a Dravidian language, has a richer agglutinative morphology.  It
encodes person, number, and gender (PNG) on verbs via suffixes, and case is
encoded directly on nominals through a system of postpositional case suffixes
(nominative, accusative, dative, genitive, locative, instrumental, etc.).  As
a result, the \texttt{Case} feature is expected to be more richly represented
in Telugu than in Hindi.

\textbf{Key expected contrasts (based on known typology):}
\begin{itemize}[leftmargin=1.5em]
  \item \texttt{Case}: More diverse values in Telugu (agglutinative case
        suffixes) than in Hindi (analytic postpositions).
  \item \texttt{Gender}: Distinct in both, but the categories differ: Hindi has
        masculine/feminine; Telugu has a three-way natural-gender distinction
        (masculine, feminine, neuter/non-human).
  \item \texttt{Number}: Present in both, with similar singular/plural contrasts.
  \item \texttt{Tense/Aspect/Mood (TAM)}: Telugu encodes tense within the verb
        morphology more transparently; Hindi TAM is expressed through auxiliary
        + participle constructions, which may appear as separate tokens.
  \item \texttt{VoiceType / VerbForm}: Both languages have verbal morphology for
        finite and non-finite forms, but Telugu's participial suffixes are highly
        productive and may generate a rich \texttt{VerbForm} distribution.
\end{itemize}

These typological distinctions are directly visible in the feature frequency
plots (Figures~\ref{fig:morph} and~\ref{fig:heatmap}): \placeholder{describe
what you observe in your actual plots here, e.g.\ ``Telugu shows a notably
higher proportion of Case tokens relative to Hindi, while Hindi shows a higher
Gender token proportion, consistent with its grammaticalised gender agreement.''}

% ══════════════════════════════════════════════════════════════════════════════
\newpage
\section{Experimental Data Analysis: Psycholinguistics (R)}
% ══════════════════════════════════════════════════════════════════════════════

\subsection*{Dataset}
The dataset is \texttt{pset\_1\_data} from MIT OpenCourseWare 9.59J (Lab in
Psycholinguistics, Spring 2017).  The dependent variable of interest is
\texttt{RTlexdec} (lexical decision reaction time, log-transformed
milliseconds), referred to as RT throughout.  All analyses are performed in
R~4.x using \texttt{ggplot2} and \texttt{dplyr}.

% ──────────────────────────────────────────────────────────────────────────────
\subsection{Q1: Distribution of RTs}
% ──────────────────────────────────────────────────────────────────────────────

\begin{figure}[H]
  \centering
  \includegraphics[width=0.75\textwidth]{results/Q1_rt_histogram_all.png}
  \caption{Histogram of RT for all subjects (60 bins).}
  \label{fig:q1}
\end{figure}

The RT distribution for all subjects is \textbf{bimodal}: it exhibits two
distinct peaks.  The RTs are not normally distributed; the histogram shows a
clear separation between two clusters of responses.  The two peaks suggest the
presence of two underlying sub-populations within the data, which Q2 clarifies.

% ──────────────────────────────────────────────────────────────────────────────
\subsection{Q2: Faceted Histograms by Age Group}
% ──────────────────────────────────────────────────────────────────────────────

\begin{figure}[H]
  \centering
  \includegraphics[width=0.9\textwidth]{results/Q2_rt_histogram_faceted.png}
  \caption{Separate RT histograms for young (left) and old (right) subjects.}
  \label{fig:q2}
\end{figure}

Separating by age group reveals that the bimodal shape in Q1 arises from
\textbf{mixing two age groups}: old subjects have systematically longer RTs than
young subjects, and each group's distribution is roughly unimodal.  The
\textbf{young subject data appear approximately normally distributed}, with a
single peak and mild right skew — consistent with a log-normal RT distribution.

% ──────────────────────────────────────────────────────────────────────────────
\subsection{Q3: Z-scores for Young Subjects}
% ──────────────────────────────────────────────────────────────────────────────

From here all analyses use only young subjects ($N = \placeholder{N}$).
Z-scores are computed as $z = (x - \bar{x}) / s$.

\subsubsection*{(a) Expected percentages under normality}

Under a standard normal distribution:
\begin{align*}
P(z > 1.96)  &\approx 2.5\% \\
P(z < -1.96) &\approx 2.5\%
\end{align*}

\subsubsection*{(b) Actual percentages}

\begin{table}[H]
\centering
\caption{Expected vs.\ actual tail percentages for young subject RTs.}
\begin{tabular}{lrr}
\toprule
\textbf{Region} & \textbf{Expected (\%)} & \textbf{Actual (\%)} \\
\midrule
$z > 1.96$  & 2.50 & \placeholder{X.XX} \\
$z < -1.96$ & 2.50 & \placeholder{X.XX} \\
\bottomrule
\end{tabular}
\end{table}

The \textbf{upper tail} ($z > 1.96$) is likely to show a greater-than-expected
proportion of observations.  This is because RT distributions tend to be
right-skewed even within a single age group: extremely slow responses (outliers)
occur more often than a perfectly normal distribution would predict, whereas
extremely fast responses are bounded from below (a participant cannot respond in
0\,ms).  This asymmetry means the positive tail contains more mass than the
theoretical 2.5\%.

\subsubsection*{(c) Words with $z > 3$}

Under normality, approximately $0.13\%$ of observations would have $z > 3$.
The actual proportion is \placeholder{X.XX\%}, which is
\placeholder{higher/lower}.  Inspection of the words with $z > 3$ reveals
\placeholder{describe what is special about them — typically low-frequency
words, function words with unusual distributions, or proper nouns that are
harder to recognise in lexical decision}.  High z-scores indicate unexpectedly
long RTs, consistent with the idea that low-frequency or phonologically
unusual words take longer to recognise.

% ──────────────────────────────────────────────────────────────────────────────
\subsection{Q4: Mean vs.\ Median}
% ──────────────────────────────────────────────────────────────────────────────

\begin{table}[H]
\centering
\caption{Mean and median comparison for RT and NounFrequency (young subjects).}
\begin{tabular}{lrrrr}
\toprule
\textbf{Variable} & \textbf{Mean} & \textbf{Median} & \textbf{Difference} & \textbf{Reason} \\
\midrule
RT           & \placeholder{M} & \placeholder{Md} & \placeholder{D} & mild right skew \\
NounFrequency & \placeholder{M} & \placeholder{Md} & \placeholder{D} & extreme right skew \\
\bottomrule
\end{tabular}
\end{table}

\begin{figure}[H]
  \centering
  \includegraphics[width=0.65\textwidth]{results/Q4_nounfreq_distribution.png}
  \caption{Distribution of NounFrequency for young subjects (dashed lines: mean in red, median in blue).}
  \label{fig:q4}
\end{figure}

The mean and median of RT are close because RT (already log-transformed in this
dataset) is approximately symmetric.  NounFrequency, by contrast, shows a very
large gap: a handful of extremely high-frequency nouns (e.g.\ function-noun
hybrids, highly common words) pull the mean far above the median.  Word
frequency in natural language follows a \textbf{Zipfian distribution} — the
mean is therefore a poor summary statistic and the median better represents the
``typical'' word.

% ──────────────────────────────────────────────────────────────────────────────
\subsection{Q5: t-test — Words Starting with `p'}
% ──────────────────────────────────────────────────────────────────────────────

A two-sided Welch t-test compared the RTs of words beginning with `p' to all
other words.

\begin{table}[H]
\centering
\caption{Two-sided t-test results: p-words vs.\ other words.}
\begin{tabular}{lr}
\toprule
\textbf{Statistic} & \textbf{Value} \\
\midrule
$t$ & \placeholder{t} \\
$df$ & \placeholder{df} \\
$p$-value & \placeholder{p} \\
Mean RT (p-words) & \placeholder{M1} \\
Mean RT (others)  & \placeholder{M2} \\
Significant at 95\%? & \placeholder{Yes/No} \\
\bottomrule
\end{tabular}
\end{table}

\placeholder{If not significant:} The difference in mean RT between words
starting with `p' and other words is not statistically significant at the 95\%
confidence level.  This suggests that the initial phoneme /p/ alone does not
reliably predict lexical decision speed, and any observed numerical difference
could be due to chance or to the confound of differing word-frequency
distributions between the two groups.

% ──────────────────────────────────────────────────────────────────────────────
\subsection{Q6: Boxplot — Noun RT vs.\ Verb RT}
% ──────────────────────────────────────────────────────────────────────────────

\begin{figure}[H]
  \centering
  \includegraphics[width=0.6\textwidth]{results/Q6_boxplot_noun_verb.png}
  \caption{Boxplot of RT for nouns (N) and verbs (V) — young subjects.}
  \label{fig:q6}
\end{figure}

Outliers appear predominantly \textbf{above} the boxes (upper whiskers), which
is consistent with the right-skewed nature of RT: extreme slow responses are
more common than extreme fast ones.  The \texttt{fivenum()} output corresponds
directly to the boxplot's five reference points: minimum, Q1, median, Q3,
maximum (after excluding beyond-whisker outliers in ggplot2's default
implementation).

\begin{table}[H]
\centering
\caption{\texttt{fivenum()} output for noun and verb RT (young subjects).}
\begin{tabular}{lrrrrr}
\toprule
\textbf{Category} & \textbf{Min} & \textbf{Q1} & \textbf{Median} & \textbf{Q3} & \textbf{Max} \\
\midrule
Nouns & \placeholder{} & \placeholder{} & \placeholder{} & \placeholder{} & \placeholder{} \\
Verbs & \placeholder{} & \placeholder{} & \placeholder{} & \placeholder{} & \placeholder{} \\
\bottomrule
\end{tabular}
\end{table}

% ──────────────────────────────────────────────────────────────────────────────
\subsection{Q7: Bar Plot with 95\% Confidence Intervals}
% ──────────────────────────────────────────────────────────────────────────────

\begin{figure}[H]
  \centering
  \includegraphics[width=0.6\textwidth]{results/Q7_barplot_ci_noun_verb.png}
  \caption{Mean RT with 95\% CI for nouns and verbs (young subjects).}
  \label{fig:q7}
\end{figure}

The 95\% confidence intervals are computed as $\bar{x} \pm 1.96 \cdot
\text{SE}$, where $\text{SE} = s / \sqrt{n}$.  \placeholder{Comment on whether
the CIs overlap.  Non-overlapping CIs suggest a statistically significant
difference; overlapping CIs suggest the difference may not reach significance.}

% ──────────────────────────────────────────────────────────────────────────────
\subsection{Q8: Boxplot by Initial Letter}
% ──────────────────────────────────────────────────────────────────────────────

\begin{figure}[H]
  \centering
  \includegraphics[width=\textwidth]{results/Q8_boxplot_by_letter.png}
  \caption{RT distribution by first letter of the word (young subjects).}
  \label{fig:q8}
\end{figure}

Some letters show notably \placeholder{higher/lower} median RTs.  Letters that
initiate many high-frequency words (e.g.\ \textit{t}, \textit{a}) tend to have
lower median RTs, while infrequent initial letters (e.g.\ \textit{x}, \textit{z})
may show higher medians due to the predominance of low-frequency words in those
bins.  The spread of each box also varies, reflecting the heterogeneity of word
frequencies for different initial letters.

% ──────────────────────────────────────────────────────────────────────────────
\subsection{Q9: Words Starting with Two Consonants}
% ──────────────────────────────────────────────────────────────────────────────

Words whose first two characters are both consonants were identified and
compared to all other words using a two-sided Welch t-test.

\begin{table}[H]
\centering
\caption{Two-sided Welch t-test: two-consonant-initial words vs.\ others.}
\begin{tabular}{lr}
\toprule
\textbf{Statistic} & \textbf{Value} \\
\midrule
$t$ & \placeholder{t} \\
$df$ & \placeholder{df} \\
$p$-value & \placeholder{p} \\
Mean RT (CC-words) & \placeholder{M1} \\
Mean RT (others)   & \placeholder{M2} \\
Significant at 95\%? & \placeholder{Yes/No} \\
\bottomrule
\end{tabular}
\end{table}

\begin{figure}[H]
  \centering
  \includegraphics[width=0.65\textwidth]{results/Q9_boxplot_two_consonants.png}
  \caption{RT for words starting with two consonants vs.\ all others.}
  \label{fig:q9}
\end{figure}

\textbf{Discussion.}
\placeholder{Fill in after running.  Typical pattern: CC-initial words tend to
have longer RTs.  This could reflect (a) lower average word frequency for
CC-onset words in English (e.g.\ `str-', `gl-' are rarer), or (b) a genuine
phonological processing cost associated with parsing a complex onset cluster,
consistent with the Obligatory Contour Principle and articulatory/perceptual
complexity.  If the difference is significant, cite the t-value and p-value.
If not, note that frequency differences between groups may be a confound that
swamps the phonological effect.}

% ══════════════════════════════════════════════════════════════════════════════
\section{Conclusion}
% ══════════════════════════════════════════════════════════════════════════════

This report has presented a systematic corpus-based comparison of Hindi and
Telugu using dependency treebanks, and an exploratory analysis of reaction time
data from a lexical decision experiment.  Key findings are:

\begin{itemize}[leftmargin=1.5em]
  \item Both Hindi and Telugu are SOV languages with right-skewed dependency
        distance distributions, consistent with the typological prediction that
        verb-final languages have longer average dependency distances than
        verb-initial ones.
  \item The Mann-Whitney U test \placeholder{confirms/does not confirm} a
        statistically significant difference in mean dependency distance,
        suggesting that within the SOV category, morphological richness and
        postpositional structure further shape syntactic locality.
  \item Morphological features reveal the expected Indo-Aryan vs.\ Dravidian
        contrast: richer agglutinative case morphology in Telugu and grammatical
        gender agreement in Hindi.
  \item In the psycholinguistic dataset, RT distributions are bimodal across
        age groups, approximately normal within the young cohort, and
        right-skewed overall.  Word frequency (NounFrequency) is strongly
        Zipfian, making the median a more appropriate central tendency measure
        than the mean.
\end{itemize}

% ══════════════════════════════════════════════════════════════════════════════
\section*{Code Repository}
% ══════════════════════════════════════════════════════════════════════════════

All code is available at: \url{\placeholder{https://github.com/USERNAME/REPO}}

\end{document}
